% Synthetisches Testdokument: KOMA-Bericht mit biblatex/Biber, PDF/A-Paket,
% Verzeichnissen im Inhaltsverzeichnis, Abkürzungen und Anhang.
\documentclass[a4paper,12pt,twoside]{scrreprt}
\usepackage[utf8]{inputenc}
\usepackage[T1]{fontenc}
\usepackage[ngerman]{babel}
\usepackage{csquotes}
\usepackage{graphicx}
\usepackage{acronym}
\usepackage[a-2b]{pdfx}
\usepackage[style=ieee,backend=biber]{biblatex}
\addbibresource{quellen.bib}

\begin{document}
\begin{titlepage}
  \centering
  \vspace*{3cm}
  {\LARGE\bfseries Titel der Beispielarbeit\par}
  \vspace{1cm}
  {\large Bachelorarbeit\par}
  \vfill
  {\large Erika Beispiel\par}
\end{titlepage}

\tableofcontents
\clearpage
\phantomsection
\addcontentsline{toc}{chapter}{Abbildungsverzeichnis}
\listoffigures

\chapter*{Abkürzungsverzeichnis}
\addcontentsline{toc}{chapter}{Abkürzungsverzeichnis}
\begin{acronym}[SPS]
  \acro{SPS}{Speicherprogrammierbare Steuerung}
\end{acronym}

\chapter{Einleitung}
Eine \ac{SPS} wird in~\parencite{muster2020} beschrieben, siehe auch \textcite{beispiel2021}.

\section{Motivation}
Text mit \enquote{Anführungszeichen}.

\appendix
\chapter{Zusatzmaterial}
Anhang.

\clearpage
\phantomsection
\addcontentsline{toc}{chapter}{Literaturverzeichnis}
\printbibliography
\end{document}
